% Source: sources/returns/S6_04_MAPS_MANIFOLDS/sections/manifolds.tex:2844–3023
\paragraph{Theorem D.2.14 (Proved: the twisted Smale suspension).}
Let $N_m$ be the unique closed simply connected spin five-manifold
with $H_2(N_m;\mathbb Z)=(\mathbb Z/m)^2$, so that $N_1=S^5$.
Define its twisted suspension by
\[
 \widetilde\Sigma_0 N_m
 =\bigl((N_m\setminus\mathring D^5)\times S^1\bigr)
       \cup_{\widetilde\alpha}(S^4\times D^2),
 \qquad \widetilde\alpha(x,z)=(\alpha_zx,z),
 \tag{D.S9}
\]
where $\alpha:S^1\to SO(5)$ represents its nontrivial
fundamental-group element. This specifies the framing in
\cite[Definition 5.1]{manGR25}; it is not replaced by an
unspecified choice between the two suspension framings.
Every simply connected member of the four families in
D.2.13 is oriented diffeomorphic to
\[
                         \widetilde\Sigma_0 N_m.
 \tag{D.S10}
\]
In particular its oriented smooth type depends only on $m$,
and these types are distinct for distinct positive $m$.
It is spin and has $p_1=0$. The proof retains the full
refined invariant required when $m$ is even; vanishing of
$p_1$ alone is not used as a classification criterion.
\begin{proof}
\emph{The angular action and its fixed curve.}
Multiplication of the line-bundle coordinate by
$e^{i\theta}$ defines the global circle action of direction
$\delta$. It commutes with the contracting quotient,
the toroidal gluings and the logarithmic translations.
On the noncompact line-bundle core the action is free:
a nontrivial contracting iterate cannot agree with a unit
scalar, by the strict norm contraction. On the four
weight-zero one-node fillings it is free for the same
reason on the normalized $\mathbb P^1\times\mathbb C^*$
annulus, and it remains free on its double curve.

The primitive hexagonal cusp is different. In the height-one
toric lattice, use the cones on the two triangles of a unit
square, with the diagonal joining $(1,0)$ to $(0,1)$.
The cocharacter $\delta=(0,1,0)$ has weights $(-1,0,1)$
in each of their smooth affine charts, up to permutation.
Every nonzero weight is a unit, so the action is semifree
there. Its fixed curves correspond to vertical edges, which
form a single orbit under the two period translations.
They therefore give one embedded $C\simeq\mathbb P^1$
in the quotient. The adjacent wall relation is
\[
 (1,0,1)+(-1,1,1)=(0,0,1)+(0,1,1).
\]
It gives the two normal degrees $-1,-1$. With the circle
weights retained, the normal bundle is
\[
 \nu_C=\mathcal O(-1)_{+1}\oplus\mathcal O(-1)_{-1}.
 \tag{D.S11}
\]
Holomorphic unit factors in the actual annular charts
commute with the angular action and deform equivariantly
to the toric factors, so these are the actual smooth
normal data. This also gives a direct chart verification
of the fixed edge described in Engel's final remarks.

On the added smooth logarithmic fibre, the lattice is
$\Lambda'=\Lambda+\mathbb Z(v/m)$. If a scalar angle
$s\delta$ lies in this lattice, applying $\chi=\gamma^*$
shows that $kn/m$ is an integer for its $v/m$ coefficient
$k$. Since here $n=\pm1$, it follows that $m\mid k$,
and then $s$ is an integer. Thus there are no new finite
stabilizers at this fibre. The global circle action is
semifree with fixed locus exactly $C$.

\emph{Equivariant surgery and the spin quotient.}
The specialization $H_2(T^4)\to H_2$ of the cusp is onto,
by the integral annulus calculation used in D.2.12.
Consequently $[C]$ maps into the cyclic fibre-coinvariant
summand generated by $s$ in (D.S7). Its $\delta$ normal
column imposes $ns=0$; hence $[C]=0$ when $n=\pm1$.
Since $X$ is simply connected, Hurewicz makes $C$
null-homotopic.

Conjugating the second complex coordinate in (D.S11)
identifies its equivariant real bundle with
$\mathcal O(-1)\oplus\mathcal O(1)$ carrying equal
scalar weights. Its transition loop is
$\operatorname{diag}(z^{-1},z)\in SU(2)$ and is
null-homotopic. Thus an equivariant tubular neighborhood
is $S^2\times D^4$, with trivial action on the first
factor and the free Hopf action on each punctured
normal sphere. Conjugating one normal coordinate reverses
normal orientation, so simultaneously reverse the model
$S^2$ orientation in this tubular identification; the
ambient orientation is then preserved.
Replace it equivariantly by $D^3\times S^3$.
The resulting $Y$ has a free circle action.

Two homotopic embeddings of $S^2$ in a six-manifold
are isotopic in this dimension; for instance apply
\cite[Theorem 2.9]{manGR25}, since $6>2\cdot2+1$.
Therefore the null-homotopic $C$ is isotopic to a small
unknotted sphere in a ball. Its framing is unique because
$\pi_2(SO(4))=0$. The underlying surgery is consequently
\[
                         Y\cong X\#(S^3\times S^3).
 \tag{D.S12}
\]
In particular $Y$ is simply connected.

The complex manifold $X$ is spin: $H^2(X;\mathbb Z)=0$
forces $c_1(X)=0$. Its spin structure is unique.
The circle action lifts with period $2\pi$ to this spin
structure. To see the period, lift the universal covering
action of $\mathbb R$ and inspect it at a point of $C$:
the parity of the isotropy weights is $1+(-1)=0$.
Thus the full $2\pi$ rotation acts trivially on the spin
lift, rather than by its nontrivial deck involution.
The equivariant trivialization of (D.S11) and its standard
Hopf replacement extend this same lift over the surgery.
The free action on $Y$ therefore has an equivariant spin
structure. Quotienting its horizontal spin frames shows
that the oriented orbit manifold $B=Y/S^1$ is spin.

The homotopy sequence gives $\pi_1(B)=0$ and the split
abelian-group sequence
\[
 0\longrightarrow H_2(Y)\longrightarrow H_2(B)
       \xrightarrow{\langle e,-\rangle}\mathbb Z
       \longrightarrow0,
\]
where $e$ is the Euler class of $Y\to B$. Thus
$H_2(B)=\mathbb Z\oplus(\mathbb Z/m)^2$ and $e$ is
primitive. The spin case of the oriented Smale--Barden
classification, with its realization of homology
isomorphisms, gives
\[
 (B,e)\cong\bigl((S^2\times S^3)\#N_m,(x,0)\bigr),
 \tag{D.S13}
\]
where the sign of the generator $x$ on the first summand
is chosen by the oriented circle-bundle convention.
The precise primary statements are
\cite[Theorems 2.2--2.3]{manBarden65}; the spin uniqueness
for $N_m$ is also recalled in \cite[Section 4.3]{manReiser24}.

\emph{The framing and cancellation.}
Apply \cite[Theorem A]{manGR25} to (D.S13), with
$B_1=S^2\times S^3$ and $B_2=N_m$.
The Euler class on $B_1$ is primitive, its circle-bundle
total space is $S^3\times S^3$, and $B_1$ is spin.
The spin branch of that theorem specifies the twisted
framing, and yields
\[
          Y\cong(S^3\times S^3)\#\widetilde\Sigma_0 N_m.
 \tag{D.S14}
\]
By \cite[Lemma 5.2]{manGR25}, the latter suspension is
simply connected and spin, has the same $H_2$ as $N_m$,
and has $b_3=0$. Zhubr's complete oriented classification
\cite[Theorem 1]{manZhubr75} permits cancellation of the
$S^3\times S^3$ summand in (D.S12)--(D.S14): connected
sum with it leaves $H_2$, the refined class $p$ and the
class $\mu$ unchanged, and adds two to $b_3$.
The diffeomorphism of the stabilized manifolds identifies
these full invariants; deleting the common increase of
$b_3$ identifies the classification data of the two
unstabilized manifolds. This proves (D.S10), including
the two-primary data. No general cancellation theorem
outside this classified class is being assumed.

Finally $TY\simeq\pi^*TB\oplus\mathbb R$, and
$H^4(B;\mathbb Z)=H_1(B;\mathbb Z)=0$, so $p_1(Y)=0$.
Restriction across (D.S12) gives $p_1(X)=0$ even though
its target $H^4(X;\mathbb Z)=(\mathbb Z/m)^2$ can be
nonzero. Distinct $m$ have distinct integral torsion
orders. For $m=1$, (D.S9) is $S^6$, also directly by
\cite[Example 5.3(1)]{manGR25}.
The argument computes the smooth type rather than
deducing it from these additive groups or from $p_1$.
\end{proof}
